% ALTERNATIVA LOCAL. No sustituye ningún PDF entregado ni el informe de raíz.
% Técnica: infografía de muestreo, con diagrama TikZ y jerarquía visual.
% Tres tiempos visuales: introducción, núcleo temático con análisis y conclusiones.
% En póster se usan bloques; los section del núcleo heredan supuestos de article.
% El reporte canónico de tres secciones conserva su fuente y PDF independientes.
\documentclass[spanish,final,xcolor={dvipsnames,table}]{beamer}
\def\itescapostertitle{Muestreo para decidir}
% Adaptador local de póster ITESCA. Requiere clase beamer; no es un article.
% Hereda el núcleo real sin copiarlo ni modificar sus configuraciones globales.
% Hook por basename de LaTeX: después de config_poster y antes de imports/cmd.
\providecommand{\itescapostertitle}{Estadística aplicada para la administración}
\providecommand{\itescaposterauthor}{Martín Jonathan de la Cruz Muñoz}
\providecommand{\itescaposterstudentid}{26130503}
\providecommand{\itescaposterdate}{25 de octubre de 2026}
\providecommand{\itescapostercourse}{Estadística aplicada para la administración}
\providecommand{\itescaposterteacher}{Dra. Yoana Elizeth Tauitmes Delgado}
\providecommand{\itescaposterwidth}{29.7}
\providecommand{\itescaposterheight}{42}
\providecommand{\itescapostercolumns}{1}
\providecommand{\itescaposterfontsize}{12}
\def\documenttitle{\itescapostertitle}
\AddToHookNext{file/config_poster.tex/after}{%
  \def\posterwidth{\itescaposterwidth}
  \def\posterheight{\itescaposterheight}
  \def\colnumbers{\itescapostercolumns}
  \def\colseparation{0.025\paperwidth}
  \def\documentfontsize{\itescaposterfontsize}
  \def\fontdocument{lmodern}
  \def\stylecitereferences{custom}
  \definecolor{itescaPosterRed}{HTML}{8F1117}
  \definecolor{itescaPosterAccent}{HTML}{D70F18}
  \definecolor{itescaPosterInk}{HTML}{181818}
  \definecolor{itescaPosterPaper}{HTML}{F4F7F7}
  \definecolor{itescaPosterTeal}{HTML}{006B67}
  \setbeamercolor{normal text}{fg=itescaPosterInk,bg=white}
  \setbeamercolor{structure}{fg=itescaPosterRed}
  \setbeamercolor{headline}{fg=itescaPosterInk,bg=white}
  \setbeamercolor{headline rule}{bg=itescaPosterAccent}
  \setbeamercolor{block title}{fg=itescaPosterRed,bg=white}
  \setbeamercolor{block body}{fg=itescaPosterInk,bg=itescaPosterPaper}
  \setbeamerfont{headline title}{size=\fontsize{24}{27},series=\bfseries}
  \setbeamerfont{headline author}{size=\fontsize{10.5}{12}}
  \setbeamerfont{headline institute}{size=\fontsize{10}{12}}
  \setbeamerfont{block title}{size=\fontsize{14}{16},series=\bfseries}
  \setbeamerfont{block body}{size=\fontsize{12}{14.2}}
  \setbeamerfont{footline}{size=\fontsize{8.5}{10}}
}
\AddToHookNext{file/imports_poster.tex/after}{%
  \usepackage{fontspec}
  \setmainfont{Arial}\setsansfont{Arial}
  \usefonttheme{professionalfonts}
  \usepackage[authoryear,round]{natbib}
  \setcitestyle{citesep={;}}
  \usepackage{tikz}
  \usetikzlibrary{arrows.meta,positioning,calc,babel}
}
\makeatletter
\def\input@path{{aulaTEX/base/Plantilla-Informe/}}
\makeatother
% Template:     Poster LaTeX
% Documento:    Núcleo del template
% Versión:      DEV
% Codificación: UTF-8
%
% Autor: Pablo Pizarro R.
%        pablo@ppizarror.com
%
% Manual template: [https://latex.ppizarror.com/poster]
% Licencia MIT:    [https://opensource.org/licenses/MIT]

% -----------------------------------------------------------------------------
% CONFIGURACIONES
% -----------------------------------------------------------------------------
% Definiciones previas
\input{src/defs}

% Ajustes usuario
\input{src/config}
\input{src/config_poster}

% -----------------------------------------------------------------------------
% IMPORTACIÓN DE LIBRERÍAS
% -----------------------------------------------------------------------------
\input{src/env/imports_poster}

% -----------------------------------------------------------------------------
% IMPORTACIÓN DE FUNCIONES Y ENTORNOS
% -----------------------------------------------------------------------------
\input{src/cmd/core}
\input{src/cmd/math}
\input{src/cmd/equation}
\input{src/cmd/image}
\input{src/cmd/title}
\input{src/cmd/other}
\input{src/cmd/column}
\input{src/cmd/poster}
\input{src/env/environments}

% -----------------------------------------------------------------------------
% IMPORTACIÓN DE ESTILOS
% -----------------------------------------------------------------------------
\input{src/style/code}
\input{src/style/other}

% -----------------------------------------------------------------------------
% CONFIGURACIÓN INICIAL DEL DOCUMENTO
% -----------------------------------------------------------------------------
\input{src/cfg/init_poster}

% -----------------------------------------------------------------------------
% CONFIGURACIÓN DE PÁGINA Y ENCABEZADOS
% -----------------------------------------------------------------------------
\input{src/cfg/page_poster}

% -----------------------------------------------------------------------------
% CONFIGURACIONES FINALES
% -----------------------------------------------------------------------------
\input{src/cfg/final_poster}

% Marco local de ancho útil explícito; evita la fila exterior de columnas vacías
% del entorno poster. Usa el mismo núcleo beamerposter, cabecera, pie y bloques.
\setbeamersize{text margin left=\sepwidth,text margin right=\sepwidth}
\newenvironment{itescaposter}[1][]{%
  \begin{frame}[fragile,t,environment=itescaposter,#1]%
}{\end{frame}}
\author{\texorpdfstring{\itescaposterauthor\quad Matrícula: \itescaposterstudentid}{\itescaposterauthor; Matrícula: \itescaposterstudentid}}
\institute{Maestría en Gestión Administrativa\\\itescapostercourse\\\itescaposterteacher}
\logoleft{width=3.4cm}{ITESCA/assets-itesca/web/logo-itesca-contorno.png}
\footercontent{\textbf{VARIANTE LOCAL DE MUESTRA}\quad GGEA01\quad\itescaposterdate}
\hypersetup{hidelinks,breaklinks=true,pdftitle={\itescapostertitle},pdfauthor={\itescaposterauthor},pdfsubject={Infografía alternativa local de muestreo; sin envío a plataforma}}
% Espaciado local de bloques. El núcleo aporta beamerposter, columnas, cabecera,
% pie, comandos matemáticos, imágenes, tablas y tamaño beamerposter.
\setbeamertemplate{block begin}{%
  \begin{beamercolorbox}[wd=\linewidth,sep=0.09cm]{block title}%
    \usebeamerfont{block title}\insertblocktitle
  \end{beamercolorbox}%
  \begin{beamercolorbox}[wd=\linewidth,sep=0.18cm]{block body}%
    \usebeamerfont{block body}\justifying
    \setlength{\parindent}{0pt}\setlength{\parskip}{0.35em}%
}
\setbeamertemplate{block end}{\end{beamercolorbox}\vspace{0.22cm}}
\newenvironment{itescaposterpanel}[1]{%
  \setbeamercolor{block title}{fg=itescaPosterTeal,bg=white}%
  \setbeamerfont{block title}{size=\fontsize{13}{15},series=\bfseries}%
  \begin{block}{#1}%
}{\end{block}}

\newlength{\itescapanelwidth}
\setlength{\itescapanelwidth}{0.485\colwidth}
\begin{document}
\templatePagecfg\templateFinalcfg
\setlength{\parindent}{0pt}
\begin{itescaposter}
\setstretch{1}\fontsize{12}{14.2}\selectfont\setlength{\parindent}{0pt}
\begin{block}{Introducción}
El muestreo selecciona unidades de una población para conocer sus características sin observarla completa. En administración reduce costo y tiempo, pero una conclusión defendible depende del diseño, la cobertura y la respuesta; reunir más datos no garantiza representar al grupo de interés \citep{bradley2021,holmesBusiness2023}.
\end{block}

\begin{block}{Muestreo: del universo a la decisión}
\centering
\begin{tikzpicture}[
 box/.style={draw=itescaPosterTeal,fill=white,rounded corners=2pt,minimum height=1.45cm,text width=7.5cm,align=center,font=\fontsize{12}{14}\selectfont},
 every path/.style={line width=1pt},>=Stealth]
\node[box] (population) {\textbf{Población}\par Clientes del periodo que interesa conocer};
\node[box,right=1.15cm of population] (frame) {\textbf{Marco muestral}\par Registro que permite seleccionarlos};
\node[box,right=1.15cm of frame] (sample) {\textbf{Muestra}\par Clientes seleccionados para observar};
\draw[->,itescaPosterTeal] (population.east)--(frame.west);
\draw[->,itescaPosterTeal] (frame.east)--(sample.west);
\end{tikzpicture}
\end{block}

\begin{columns}[t,totalwidth=\colwidth]
\begin{column}{\itescapanelwidth}
\begin{itescaposterpanel}{01 / Diseños probabilísticos}
Cada unidad tiene probabilidad conocida y positiva de inclusión; no siempre es igual.

\textbf{Aleatorio simple.} Todas las muestras del mismo tamaño tienen la misma probabilidad: sortear clientes de un registro completo.

\textbf{Estratificado.} Seleccionar al azar dentro de todos los estratos: clientes por canal de compra.

\textbf{Sistemático.} Inicio aleatorio y cada $k$ unidades: revisar cada décima factura; evitar periodicidades.

\textbf{Conglomerados.} Sortear grupos, como sucursales, y estudiar sus unidades. Reduce desplazamientos, pero la semejanza interna puede disminuir precisión \citep{holmesBusiness2023}.
\end{itescaposterpanel}

\begin{itescaposterpanel}{02 / ¿Cuántas respuestas?}
\textbf{Ejemplo hipotético:} proporción de clientes satisfechos; muestra aleatoria simple, aproximación normal y población grande respecto de la muestra.
\[
 n_0=\frac{z_{\alpha/2}^{\,2}p(1-p)}{e^2}
 =\frac{1.96^2(0.5)(0.5)}{0.05^2}=384.16
\]
\textbf{Planificar 385 respuestas completas.} Confianza del 95\%, margen de cinco puntos porcentuales y $p=0.5$, que maximiza la variabilidad de planificación \citep{holmesBusiness2023}.

Con respuesta prevista del 80\%:
\[
 \lceil385/0.80\rceil=482\quad\text{invitaciones}.
\]
No garantiza esa respuesta ni corrige sesgo. Otros diseños requieren ajustes. Para contrastar hipótesis, también se especifican potencia y efecto mínimo detectable \citep{nistHandbook}.
\end{itescaposterpanel}
\end{column}
\begin{column}{\itescapanelwidth}
\begin{itescaposterpanel}{03 / Selección no probabilística}
No permite conocer la probabilidad de inclusión de cada unidad.

\textbf{Conveniencia.} Entrevistar a quienes están disponibles, como compradores presentes en una tienda.

\textbf{Autoselección.} Responder por iniciativa propia, como en una encuesta abierta en redes sociales.

Sirve para explorar; no atribuirle el margen de error de una muestra aleatoria sin supuestos adicionales. \textbf{Más respuestas no corrigen por sí solas el sesgo.} Una selección aleatoria dentro de un marco incompleto tampoco cubre automáticamente a toda la población \citep{bradley2021}.
\end{itescaposterpanel}

\begin{itescaposterpanel}{04 / Aplicaciones administrativas}
\textbf{Calidad de facturación.} Seleccionar facturas con inicio aleatorio, estimar la proporción con errores y decidir dónde auditar. Los errores detectados no demuestran sus causas.

\textbf{Satisfacción por canal.} Estratificar clientes de tienda y venta digital, seleccionar al azar en ambos y ponderar según el diseño. Escuchar solo a seguidores digitales excluye otros compradores.

La comparación conduce a una conclusión práctica: \textbf{definir población, marco y selección antes de ampliar la encuesta}. El porcentaje orienta decisiones cuando se conocen tanto las unidades observadas como las excluidas \citep{bradley2021,holmesBusiness2023}.
\end{itescaposterpanel}
\end{column}
\end{columns}

\begin{block}{Conclusiones}
Una decisión responsable revisa cobertura, selección y no respuesta antes de celebrar el tamaño de la base. Conviene invertir primero en datos pertinentes y después en cantidad: así el ahorro del muestreo no se convierte en un costo por decisiones sesgadas. Este criterio permite delimitar qué se puede inferir y qué evidencia adicional necesita la gerencia.
\end{block}

\begin{block}{Referencias}
\fontsize{9.5}{11.1}\selectfont\setstretch{1}\raggedright
\urlstyle{same}\Urlmuskip=0mu plus 1mu
\setlength{\emergencystretch}{2em}
\renewcommand{\bibsection}{}\setlength{\bibsep}{2pt}\setlength{\bibhang}{0.65cm}
\begin{thebibliography}{3}
\bibitem[Bradley et al.(2021)]{bradley2021}
Bradley, V. C., Kuriwaki, S., Isakov, M., Sejdinovic, D., Meng, X.-L., \& Flaxman, S. (2021). Unrepresentative big surveys significantly overestimated US vaccine uptake. \textit{Nature, 600}(7890), 695--700. \url{https://doi.org/10.1038/s41586-021-04198-4}
\bibitem[Holmes et al.(2023)]{holmesBusiness2023}
Holmes, A., Illowsky, B., \& Dean, S. (2023). \textit{Introductory business statistics} (2.\textsuperscript{a} ed.). OpenStax. \url{https://openstax.org/details/books/introductory-business-statistics-2e}
\bibitem[NIST y SEMATECH(2003)]{nistHandbook}
National Institute of Standards and Technology \& SEMATECH. (2003). \textit{NIST/SEMATECH e-Handbook of statistical methods}. \url{https://doi.org/10.18434/M32189}
\end{thebibliography}
\end{block}
\end{itescaposter}
\end{document}
